% =====================================================================
%  Johan David Rodriguez Castro — Curriculum Vitae (cuerpo compartido)
%
%  No compilar este archivo directamente. Usar:
%    johan-rodriguez-cv-onepage.tex   → 1 página, para postular (ATS)
%    johan-rodriguez-cv.tex           → 2 páginas, versión detallada
%
%  Reglas de formato, tomadas de la práctica estándar para perfiles sin
%  experiencia laboral larga y verificadas con `pdftotext` sobre este PDF:
%    · Una sola columna. Nada de cajas de texto ni columnas laterales.
%    · Ninguna tabla de dos columnas para las skills: al extraer el texto,
%      un ATS lee todas las etiquetas primero y después todos los valores,
%      y la asociación etiqueta→keywords se pierde. Se usan líneas simples.
%    · Todo es texto real. Nada dentro de una imagen.
%    · Los proyectos se escriben con la misma estructura que un empleo:
%      título, entidad, fechas y bullets con cifras.
% =====================================================================
\providecommand{\cvvariant}{full}
\def\cvonepagename{onepage}
\ifx\cvvariant\cvonepagename \def\cvfontsize{9pt}\else \def\cvfontsize{10pt}\fi

% extarticle (paquete extsizes) admite 9pt; la clase article estándar no.
\documentclass[\cvfontsize,a4paper]{extarticle}

\newif\ifonepage
\ifx\cvvariant\cvonepagename \onepagetrue \else \onepagefalse \fi

\usepackage[T1]{fontenc}
\usepackage[utf8]{inputenc}
\usepackage{tgtermes}                 % cuerpo: compatible con Times
\usepackage{tgheros}                  % títulos: compatible con Helvetica
\usepackage{microtype}

\ifonepage
  \usepackage[a4paper, top=0.85cm, bottom=0.85cm, left=1.2cm, right=1.2cm]{geometry}
\else
  \usepackage[a4paper, top=1.2cm, bottom=1.2cm, left=1.4cm, right=1.4cm]{geometry}
\fi

\usepackage{xcolor}
\definecolor{accent}{HTML}{1F4E79}
\definecolor{subtle}{HTML}{555555}

\usepackage[
  colorlinks=true, urlcolor=accent, linkcolor=accent,
  pdftitle={Johan David Rodriguez Castro — CV},
  pdfauthor={Johan David Rodriguez Castro},
  pdfsubject={Security \& AI Engineer},
  pdfkeywords={security engineering, detection engineering, supply-chain security,
               vulnerability management, applied machine learning, Python, AWS, Terraform}
]{hyperref}

\usepackage{titlesec}
\usepackage{enumitem}
\usepackage{ragged2e}

\titleformat{\section}[block]
  {\sffamily\bfseries\color{accent}\fontsize{10.5}{12}\selectfont}
  {}{0em}{\MakeUppercase}[\vspace{-0.9ex}{\color{accent}\rule{\linewidth}{0.7pt}}]
\titlespacing*{\section}{0pt}{0.9ex plus 0.2ex}{0.6ex}

\setlist[itemize]{
  leftmargin=1.1em, topsep=0.15ex, itemsep=0.05ex, parsep=0pt,
  label=\textcolor{accent}{\textbullet}
}

\setlength{\parindent}{0pt}
\pagestyle{empty}

% ── Entrada de empleo o proyecto ─────────────────────────────────────
% \entry{título}{subtítulo}{fechas · entidad}
% Se usa \hfill dentro de un solo párrafo (no una tabla) para que la
% extracción de texto devuelva título y fecha en la MISMA línea.
\newcommand{\entry}[3]{%
  \vspace{0.5ex}\par\noindent
  \textbf{#1}\ifx\relax#2\relax\else{\,---\,\itshape #2}\fi
  \hfill{\footnotesize\color{subtle}#3}\par\nobreak\vspace{-0.1ex}%
}

% Línea de enlaces bajo una entrada.
\newcommand{\evidence}[1]{{\footnotesize\color{subtle}#1}\par\nobreak}

% Entrada con evidencia: \project{título}{subtítulo}{enlaces}{fechas · entidad}
% En 1 página los enlaces se pliegan en la línea del título para no gastar
% una línea por proyecto.
\newcommand{\project}[4]{%
  \ifonepage
    \entry{#1}{#2}{#3 \ {\color{subtle}$\vert$} \ #4}%
  \else
    \entry{#1}{#2}{#4}\evidence{#3}%
  \fi
}

% Fila de skills: una línea de texto corrido, no una celda de tabla.
\newcommand{\skill}[2]{%
  \par\vspace{0.35ex}\noindent\textbf{#1:}\ #2\par
}

\begin{document}

% =====================================================================
\begin{center}
  {\sffamily\bfseries\color{accent}\fontsize{18}{21}\selectfont
   JOHAN DAVID RODRIGUEZ CASTRO}\\[0.5ex]
  {\footnotesize\color{subtle}
   Bucaramanga, Santander, Colombia
   \ $\vert$ \ +57 313 878 4948
   \ $\vert$ \ \href{mailto:johan.rc2020@gmail.com}{johan.rc2020@gmail.com}}\\[0.4ex]
  {\footnotesize
   \href{https://github.com/Yoyagm}{github.com/Yoyagm}
   \ {\color{subtle}$\vert$} \
   \href{https://www.linkedin.com/in/johan-rodriguez-97bb29323/}{linkedin.com/in/johan-rodriguez-97bb29323}
   \ {\color{subtle}$\vert$} \
   \href{https://johan-rodriguez.vercel.app}{johan-rodriguez.vercel.app}}
\end{center}
\vspace{0.2ex}

% =====================================================================
\section{Professional Profile}
\justifying
\ifonepage
Final-semester Systems and Software Engineering student (UPB Bucaramanga) working at the intersection of
security engineering and applied machine learning. I ship tools rather than findings: an open-source
supply-chain scanner, a phishing classifier that audited its own benchmark and reported the honest number
instead of the flattering one, and a vulnerability prioritisation dashboard that measured away the
advantage it was built to show. Verified C2 English (EF SET 80/100).
\else
Final-semester Systems and Software Engineering student (UPB Bucaramanga) working at the intersection
of security engineering and applied machine learning. I ship tools rather than findings: an open-source
supply-chain scanner, a phishing classifier that audited its own benchmark and published the honest
number instead of the flattering one, and a vulnerability prioritisation dashboard that measured away
the advantage it was built to demonstrate. Equally comfortable in cloud infrastructure --- AWS,
Terraform, CI/CD --- and full-stack delivery, with Scrum leadership experience and verified C2 English
(EF SET 80/100). Seeking a Security, Detection Engineering or Data \& AI role.
\fi

% =====================================================================
\section{Technical Skills}
\vspace{0.2ex}
\skill{Security \& Detection}{Software supply-chain security (slopsquatting), detection-as-code,
  MITRE ATT\&CK mapping, risk-based vulnerability prioritisation (CISA KEV / EPSS / CVSS),
  phishing and URL analysis, TOTP 2FA, secure SDLC and CI security gates}
\skill{Data, ML \& AI}{LightGBM, XGBoost, scikit-learn; data-leakage detection and honest evaluation
  design; SHAP explainability; grouped, temporal and walk-forward validation; probability calibration;
  FinBERT NLP; CrewAI multi-agent orchestration; RAG with ChromaDB; pandas / NumPy}
\skill{Cloud \& DevOps}{AWS (IAM, VPC, EC2, Auto Scaling, ALB, RDS Multi-AZ, ECR, CloudWatch),
  Terraform (IaC), Docker, GitHub Actions CI/CD, k6, Google Cloud (BigQuery ML, Cloud Storage, DLP)}
\skill{Programming}{Python, TypeScript / JavaScript, SQL, Dart, HCL, LaTeX}
\skill{Databases}{PostgreSQL, MySQL, SQLite, DuckDB, ChromaDB; hybrid SQL / NoSQL architectures}
\skill{Methods \& Tools}{Agile / Scrum, spec-driven development, pytest, Playwright, mypy strict,
  UML modelling, Git}

% =====================================================================
\section{Experience}

\project{Developer / Technical Support}{}%
  {GOATGuard: \href{https://github.com/Yoyagm/goatguard-app}{github.com/Yoyagm/goatguard-app}}%
  {Jan 2026 -- Present \ $\vert$ \ Pontificia Universidad Bolivariana}
\begin{itemize}
  \item Designed and built GOATGuard, a full-stack network security monitoring platform
        (FastAPI + PostgreSQL + Flutter), to close mobility gaps for administrators managing
        distributed institutional assets.
  \item Built an automatic classification engine for network assets (routers, printers, cameras) with
        real-time monitoring and KPI dashboards, and proactive detection of port scans and intrusion
        attempts, applying Blue Team and detection engineering principles.
  \ifonepage\else
  \item Hardened the mobile client with a full TOTP 2FA flow (QR enrolment, backup and recovery codes)
        and JWT held in OS-level secure storage behind a global 401 interceptor.
  \item Provide technical support to end users across the institution.
  \fi
\end{itemize}

\entry{Scrum Master}{}{Jun 2025 -- Nov 2025 \ $\vert$ \ Pontificia Universidad Bolivariana}
\begin{itemize}
  \item Managed the software development lifecycle under Scrum, facilitating ceremonies and optimising
        the Product Backlog to ensure incremental value delivery.
  \item Contributed to the design of a high-availability infrastructure targeting 100{,}000 concurrent
        users with a sub-500\,ms latency objective, on hybrid SQL / NoSQL / in-memory data architectures.
\end{itemize}

\ifonepage\else
\entry{Assistant, Sports Events}{}{May 2025 -- Present \ $\vert$ \ EM Tournaments and Competitions}
\begin{itemize}
  \item Support the organisation and execution of sports tournaments and competitive events.
\end{itemize}
\fi

% =====================================================================
\section{Key Technical Projects}

\project{SlopGuard}{Open-Source Supply-Chain Security Scanner}%
  {\href{https://github.com/Yoyagm/slopguard}{github.com/Yoyagm/slopguard}}%
  {2026 \ $\vert$ \ Independent}
\begin{itemize}
  \item Built a pre-install guard that scores PyPI and npm dependencies against hallucinated,
        typosquatted and malicious names \emph{before} any code executes, using a five-layer detection
        engine with zero runtime dependencies.
  \ifonepage\else
  \item Designed a provable anti-false-positive invariant: the opt-in LLM hallucination layer is
        structurally unable to block a legitimate package on its own.
  \fi
  \item Enforced correctness with a CI gate at $\geq$90\,\% global coverage ($\geq$95\,\% on critical
        paths) and eight import-linter architecture contracts, and shipped four integration points from
        one core: CLI, pre-commit hook, GitHub Action and self-hostable SaaS.
\end{itemize}

\ifonepage
\project{Phishing URL Detector}{Explainable Detection \& Leakage Audit}%
  {\href{https://yoyagm.github.io/phishing-url-detector/}{yoyagm.github.io/phishing-url-detector}}%
  {2026 \ $\vert$ \ Independent}
\else
\project{Phishing URL Detector}{Explainable Detection \& Dataset Leakage Audit}%
  {\href{https://yoyagm.github.io/phishing-url-detector/}{yoyagm.github.io/phishing-url-detector}
   \ $\vert$ \ \href{https://github.com/Yoyagm/phishing-url-detector}{github.com/Yoyagm/phishing-url-detector}}%
  {2026 \ $\vert$ \ Independent}
\fi
\begin{itemize}
  \item Audited PhiUSIIL (UCI id 967, 235{,}795 URLs) before training and measured two independent
        label leaks: a one-line rule scores 99.67\,\% accuracy, and hand-rolled lexical features on the
        raw URL --- the obvious fix --- still score 99.51\,\% while measuring nothing.
  \item Retrained on canonicalised host names only and published the honest result: ROC-AUC 0.937 and
        70.5\,\% recall at a 1\,\% false-positive rate, against phishing submitted two years after the
        training data was collected.
  \ifonepage
  \item Shipped it: exact TreeSHAP reimplemented in dependency-free JavaScript (parity-tested at 1e-12
        against the Python reference), 113 tests including leakage guards with negative controls, and a
        containerised FastAPI service.
  \else
  \item Reimplemented exact path-dependent TreeSHAP in dependency-free JavaScript, parity-tested at
        1e-12 against the Python reference over 28{,}917 attributions, so the live demo explains every
        verdict with no backend.
  \item Delivered 113 tests --- including leakage guards that ship with planted-leak negative controls
        --- a containerised FastAPI service and a nine-page technical report in LaTeX.
  \fi
\end{itemize}

\project{Vulnerability Prioritisation Dashboard}{CISA KEV + EPSS + NVD}%
  {\href{https://github.com/Yoyagm/vuln-prioritization-dashboard}{github.com/Yoyagm/vuln-prioritization-dashboard}}%
  {2026 \ $\vert$ \ Independent}
\begin{itemize}
  \item Built a reproducible ingestion pipeline over three public feeds ranking 359{,}399 live CVEs,
        with committed data contracts and schema tests that fail the build when a feed changes shape.
  \ifonepage
  \item Reproduced the industry claim that EPSS is 6.1$\times$ more efficient than CVSS, then showed
        with frozen scores over five cutoffs that most of it is evaluation artefact: EPSS wins the head
        of the queue (1.2--1.8$\times$) and loses the long tail (0.5--0.7$\times$).
  \else
  \item Reproduced the industry claim that EPSS is 6.1$\times$ more efficient than CVSS, then showed ---
        with scores frozen at five past cutoffs and only later-confirmed exploitation counted --- that
        most of it is evaluation artefact: EPSS wins the head of the queue (1.2--1.8$\times$) and loses
        the long tail (0.5--0.7$\times$).
  \fi
  \ifonepage\else
  \item Bootstrapped 95\,\% confidence intervals over the positives and guarded the tie-averaging in the
        core metric with a test that ships with a negative control; 76 tests and five ADRs.
  \fi
\end{itemize}

\ifonepage
\vspace{0.6ex}
{\footnotesize\noindent\textbf{Also:} AWS three-tier HA infrastructure in Terraform + GitHub Actions
(\href{https://github.com/Yoyagm/eval-final-dgiti}{eval-final-dgiti}) \ $\vert$ \ autonomous ML and
multi-agent market pipeline, 505 tests \ $\vert$ \ GreenPulse offline-first platform. Case studies:
\href{https://johan-rodriguez.vercel.app/en/projects}{johan-rodriguez.vercel.app/projects}.\par}
\else

\entry{AWS High-Availability Infrastructure}{Terraform, FastAPI \& CI/CD}{2026 \ $\vert$ \ UPB}
\evidence{\href{https://github.com/Yoyagm/eval-final-dgiti}{github.com/Yoyagm/eval-final-dgiti}}
\begin{itemize}
  \item Codified a three-tier architecture across two availability zones entirely in Terraform:
        VPC, ALB, Auto Scaling Group, RDS PostgreSQL Multi-AZ, ECR and CloudWatch.
  \item Enforced least privilege with a strict security-group chain --- the database has no public
        access and no egress, and is reachable only from the application tier.
  \item Automated deployment with GitHub Actions (lint, test, build, push to ECR, ASG instance refresh),
        keeping Terraform out of the hot path, and validated load balancing under k6 load against
        CloudWatch metrics.
\end{itemize}

\entry{Autonomous Market Intelligence System}{Applied ML \& Multi-Agent AI}{2025 -- Present \ $\vert$ \ Independent}
\evidence{Private repository --- walkthrough available on request}
\begin{itemize}
  \item Built an event-driven pipeline integrating market-data APIs, technical-indicator computation and
        a FinBERT sentiment model, orchestrating four CrewAI agents into daily automated briefings.
  \item Engineered a governance layer of seven risk checks plus a portfolio drawdown circuit breaker;
        medium-confidence signals escalate to human approval over Telegram and expire unanswered.
  \item Trained an XGBoost / RandomForest / LogisticRegression ensemble with walk-forward validation, a
        purged 21-day gap and Platt calibration; validated by 505 automated tests across 17 suites.
  \item Runs on Alpaca paper trading and deliberately stays there: a GO/NO-GO gate blocks live capital
        until the model clears 55\,\% accuracy (currently 52.9\,\%).
\end{itemize}

\entry{GreenPulse}{Offline-First Agronomic Monitoring Platform}{2025 \ $\vert$ \ Academic Capstone}
\begin{itemize}
  \item Migrated and redesigned the backend to FastAPI with an offline-first SQLite sync layer and a
        35-species plant catalogue including QR-code based scanning; documented in IEEE format.
\end{itemize}
\fi

% =====================================================================
\section{Education}

\entry{B.S. in Systems and Software Engineering}{}{Feb 2023 -- Present \ $\vert$ \ Pontificia Universidad Bolivariana}
\begin{itemize}
  \ifonepage
  \item Final semester. Coursework in software architecture and system design; AWS Academy Cloud
        Foundations labs 1--6 (IAM, VPC, EC2, EBS, Auto Scaling, Load Balancing).
  \else
  \item Final-semester student. Advanced coursework in software architecture, system design and
        industry best practices.
  \item AWS Academy Cloud Foundations coursework --- labs 1--6 completed: IAM, VPC, EC2, EBS,
        Auto Scaling and Load Balancing, with Terraform IaC implementations (2026).
  \fi
\end{itemize}

% =====================================================================
\section{Certifications}

\ifonepage
{\footnotesize\noindent
\textbf{Google Cloud} (2026) --- 7 skill badges: Secure Lakehouse Data \textbullet\ Protect Sensitive
Data with Data Loss Prevention \textbullet\ Create ML Models with BigQuery ML \textbullet\ Use APIs to
Work with Cloud Storage \textbullet\ Gemini Enterprise (multi-agent workflows, first application)
\textbullet\ Vibe Coding and MCP. \ \textbf{Google} (2026) --- DeepMind: Train A Small Language Model.\par
\vspace{0.2ex}\noindent
\textbf{Cisco Networking Academy} (2026) --- Introduction to Cybersecurity. \
\textbf{Hugging Face} (2026) --- AI Agents Course. \
\textbf{freeCodeCamp} (2026) --- Machine Learning with Python \textbullet\ Data Analysis with Python
\textbullet\ Scientific Computing with Python \textbullet\ Data Visualization \textbullet\ JavaScript
Algorithms and Data Structures.\par
\vspace{0.2ex}\noindent
\textbf{Languages:} Spanish --- Native \textbullet\ English --- C2 Proficient
(\href{https://cert.efset.org/es/qpqdd3}{EF SET 80/100, verified}).
{\color{subtle}All credentials publicly verifiable at
\href{https://www.linkedin.com/in/johan-rodriguez-97bb29323/details/certifications/}{linkedin.com/in/johan-rodriguez-97bb29323}.}\par}
\else
{\footnotesize\color{subtle}All credentials are publicly verifiable; full list with links at
\href{https://www.linkedin.com/in/johan-rodriguez-97bb29323/details/certifications/}{linkedin.com/in/johan-rodriguez-97bb29323}.}
\vspace{0.4ex}
\begin{itemize}
  \item \textbf{Google Cloud} (2026) --- Secure Lakehouse Data \textbullet\ Protect Sensitive Data with
        Data Loss Prevention \textbullet\ Create ML Models with BigQuery ML \textbullet\ Use APIs to Work
        with Cloud Storage \textbullet\ Orchestrate Multi-agent Workflows with Gemini Enterprise
        \textbullet\ Create Your First Gemini Enterprise Application \textbullet\ Build a Smart Cloud
        Application with Vibe Coding and MCP \textbullet\ Google DeepMind: Train A Small Language Model
  \item \textbf{Cisco Networking Academy} (2026) --- Introduction to Cybersecurity
  \item \textbf{Hugging Face} (2026) --- AI Agents Course
  \item \textbf{freeCodeCamp} (2026) --- Machine Learning with Python \textbullet\ Data Analysis with
        Python \textbullet\ Scientific Computing with Python \textbullet\ Data Visualization
        \textbullet\ JavaScript Algorithms and Data Structures
  \item \textbf{EF SET} (2026) --- English Certificate, CEFR C2 (Proficient), overall 80/100
        (reading 84, listening 89)
\end{itemize}
\fi

% =====================================================================
\ifonepage\else
\section{Languages}
Spanish --- Native \ \ {\color{subtle}$\vert$} \ \ English --- C2 Proficient
(\href{https://cert.efset.org/es/qpqdd3}{EF SET 80/100, verified})
\fi

\end{document}
